% !TEX root = ../my-thesis.tex
%
\chapter{Discussion}
\label{sec:discussion}

\section{The Added Value of the Formal Layer}
\label{sec:discussion:value}

How much the formal layer adds varies strongly across the eight topics of
Section~\ref{sec:evaluation:crosstopic}. The clearest payoff appears on
Abortion, where four preferred extensions over an undecided core of nine
arguments make competing internally consistent positions explicit, and on
Death Penalty with three extensions over a core of eight. On Gun Control,
Marijuana Legalization, and School Uniforms, by contrast, the undecided core
is empty, the grounded extension already decides every argument, and the
enumeration adds nothing beyond the cautious core. That every extracted
framework is cyclic is a precondition for this payoff rather than a merit in
itself, since in an acyclic framework the grounded, preferred, and stable
semantics coincide on a single extension~\cite{dung1995}. An honest reading is
therefore that the formal layer is most useful where the debate genuinely
splits into competing coherent camps, and closer to redundant where one side
dominates the sampled arguments.

The ablation sharpens this picture. The formal richness that the pipeline
reports is not a free observation about the debates but a consequence of one
deliberate modelling decision, the symmetrization of rebuttals, which only
typed extraction makes available. This is a feature and a caveat at once. It is a feature
because the mutual attack is the semantically correct rendering of a genuine
contradiction between conclusions. It is a caveat because every mislabelled
rebuttal buys an unjustified cycle, and Section~\ref{sec:evaluation:errors}
shows that such mislabelling occurs. The computational price of the richness
stays low. The undecided cores remain at nine or fewer arguments, and all three
semantics solve in under half a second per topic.

Read against the motivation of Chapter~\ref{sec:intro}, the multiple
extensions are more than a structural statistic. A standalone LLM that is
asked to summarise the positions of a debate returns another fluent generation
without a guarantee of internal consistency. The extensions computed here
carry that guarantee by construction, since every preferred extension is
conflict-free and admissible, and every acceptance is traceable to the
extracted attack edges that produce it. On Abortion the difference becomes
concrete, since the four extensions separate a legally framed camp from a
morally framed one, a finer distinction than the binary stance labels of the
corpus. The formal layer thus delivers what a raw model answer leaves
implicit, a set of debate positions whose joint acceptability is checkable
rather than asserted.

These results also feed back into the related work of
Chapter~\ref{sec:related}. They supply the empirical layer that
ARGUS~\cite{argus} leaves open, showing that such a pipeline produces sparse
and cyclic frameworks with small undecided cores rather than intractable
graphs. They refine the finding of Gorur et al.~\cite{gorur2025relMining} that
LLMs identify attack relations reliably, since the reference evaluation locates the
benefit of typed prompting in precision rather than recall. And they qualify
the quantitative route of ArgLLMs~\cite{argLLMs}, because a single strength
value per argument could not have represented the four incompatible Abortion
positions that the set-based semantics make explicit.

\section{Limitations}
\label{sec:discussion:threats}

Four limitations qualify the results, namely the sensitivity of the extraction to
prompt wording, the small samples, the possible familiarity of the model with
the corpus, and the reliance on a single model and a single annotator.

\paragraph{Prompt sensitivity.} The extraction is sensitive to the exact
wording of the prompts. When the attack prompt was revised during development,
among other things to add the descriptive-versus-normative rule, the accepted
edge sets of otherwise identical runs changed substantially, on some topics by
dozens of edges, and the number of extensions changed with them. The final
revision before the reference evaluation run, which tightened all three type definitions,
roughly halved the accepted typed edges on Nuclear Energy from 43 to 23. The
cross-topic and ablation numbers therefore come from one canonical batch with
one versioned prompt, the reference scores from one dedicated run with the final
revision, and all of them should be read as results about the pipeline with
these exact prompts, not
about typed prompting in the abstract. The prompt versioning of
Section~\ref{sec:method:logging} makes this dependency explicit and traceable
rather than removing it.

\paragraph{Sampling and scale.} The UKP samples are small, with seven pro and seven con
arguments per topic, which limits how far the cross-topic numbers generalise to
larger argument sets. The residual nondeterminism of the atomization, described
in Section~\ref{sec:evaluation:setup}, is controlled by the frozen atom sets
within this evaluation, but it means that an independent reproduction will
operate on a slightly different atom set.

\paragraph{Data contamination.} The UKP corpus has been publicly available
since 2018 and may therefore have been part of the training data of
GPT-OSS-120B, whose exact composition is not disclosed. In that case the model
may recognise the arguments and their stances rather than judge them from
scratch. The pairwise attack judgments, however, cannot be memorised. The reference
annotation of Section~\ref{sec:evaluation:gold} was created for this thesis on
a frozen atom set and is unpublished, and the atomized claims the model
judges mostly differ in wording from the corpus sentences. Familiarity with the
arguments could still make the task easier than it would be on genuinely
unseen debates, so the absolute scores should be read with this reservation.

\paragraph{Model and self-evaluation.} All results stem from a single model at
temperature zero. In the LLM-generated setting the arguments, the attacks, and
the confidence scores all come from the same model, so there is no independent
quality signal, and that setting is treated as a qualitative check rather than
a second evaluation. The reference annotation of Section~\ref{sec:evaluation:gold} is
created by a single annotator, the author, and no inter-annotator agreement
can be reported, a limitation that is common for a thesis but should be kept in
mind when reading the scores. The same person also designed the pipeline and
its prompts, so annotation decisions and extraction behaviour could share blind
spots, which the written guidelines and the strict object reading mitigate but
cannot rule out. The reference set is moreover small, with 16 directed edges on
a single topic, so every individual edge moves recall by six points and small
F1 differences between configurations should not be over-interpreted.

\section{Design Trade-offs}
\label{sec:discussion:tradeoffs}

Four design decisions shape the frameworks and deserve explicit discussion.

The symmetrization of rebuttals is well motivated by the absence of a
preference ordering in Dung's abstract framework, but as
Section~\ref{sec:discussion:value} discusses, it is also the single source of
every multi-extension result, which makes a mislabelled rebuttal costly. On
the manually annotated topic the mechanism proves its worth, since three of the four reverse
edges it adds are confirmed by the reference annotation. A natural refinement would be to symmetrize only
above a confidence threshold, or to let the verifier confirm specifically the
mutuality of the contradiction rather than the attack as such.

Dropping the attack type when converting to binary attack facts is consistent
with the set-based question Dung semantics answer, yet it means the pipeline
extracts information it then discards. The types, confidence scores, and
justifications are retained as metadata, and this thesis itself demonstrates
their analytical value, since the error analysis and the ablation both rely on
them. A formalism that consumes them in the semantics, such as ASPIC+ or a
weighted framework, is the obvious next step and is discussed as future work.

The confidence thresholds are design choices rather than learned values, since
the score is a routing signal and not a calibrated probability. The logged
decisions of Section~\ref{sec:evaluation:crosstopic} show that the
verification stage acts almost
exclusively as an additional reject filter, since it confirmed only six of the
68 borderline candidates across the canonical batch, and the reference evaluation run of
Section~\ref{sec:evaluation:gold} gives a first, small-sample indication that
this filtering points in the right direction. Whether the six rescued accepts
justify the extra LLM call per borderline case is a question these few cases
cannot settle. Now that a reference annotation exists, a natural follow-up
would be to tune both thresholds against it, for instance to maximise the F1
score, which was ruled out at design time because no labels were available
and tuning on the evaluation data would have leaked information.

The exhaustive pairwise querying, finally, bounds the scale the pipeline can
handle. The number of queried pairs grows quadratically with the atom count.
The 21 Nuclear Energy atoms already produce 201 pairs, and extraction takes
between two and eight minutes per topic at this size, so a realistic debate with
hundreds of arguments would produce tens of thousands of pairs and hours of
extraction. Applying the pipeline at that scale requires a cheaper candidate
filter before the LLM call, for example an embedding-based prescreening that
discards clearly unrelated pairs, at the price of missing conflicts the
filter does not anticipate.
